\begin{helper}[Situation 7.3: replace a base-extra color]
Let $G$ be a finite loopless multigraph with a finite individual edge
type, let $[k]$ be its palette, and fix a distinguished edge $e_0$, an
input coloring, an exact ordered tree-sequence $T$, a generating coloring
$\varphi$, an ETT $\mathcal E$ on these data, a
\noderef{HierarchyBase} $\mathcal B$, a
\noderef{HierarchyCounterexample} $C$, and a
\noderef{HierarchyMinimumCounterexampleContext} $\mathit{ctx}$.
Thus the context supplies properness, exactly-one-uncoloredness, all six
technical assertions at strictly lower positive rungs, fewer-divider
induction at the same rung, rung-local lexicographic minimality, and the
terminal-ladder bound $|\overline\varphi(T_n)|\ge 2n+11$.
Write $n=\mathcal E.\mathit{rounds}$, $H=C.\mathit{hierarchy}$,
$q=H.\mathit{dividerCount}$, and set
\[
 A=\noderef{HierarchyStarredStage}(H,q),\quad
 \ell=\noderef{HierarchyFullTailLength}(C),\quad
 T(y_i)=\noderef{HierarchyFullTailPrefix}(C,i),\quad
 B_{\rm pen}=T(y_{\ell-1}),\quad y_\ell=T.\mathit{vertex}(\mathit{last}).
\]
Let $T_n$ be the terminal ladder tree
$\mathcal E.\mathit{series.data.tree}((\operatorname{Fin.last}n).\mathit{castSucc})$,
and put
\[
 D=\noderef{HierarchyStoredColorsAt}(H,q),\quad
 \Gamma=\bigcup_{\rho\in D}C.\mathit{good}.\Gamma^q_\rho,\quad
 M_0=\bigcup_{v\in V(\mathcal B.\mathit{core})\setminus V(T_n)}
             \noderef{MissingColors}(G,\varphi,v).
\]
Here $V(U)=\noderef{TreeSequenceVertexSet}(U)$ and
$\overline\varphi(U)=\noderef{TreeMissingColors}(\varphi,U)$.
Assume $2\le j<\ell$, where
$j=\noderef{HierarchyPathNumber}(H)$.
Let $\alpha$ be missing on $A$ and at $y_\ell$.
Assume $\alpha\in M_0\setminus\Gamma$.
Then there are an original-relative
\noderef{HierarchyReplacementPresentation} $P$ and a color $\beta$
missing under $P.\mathit{coloring}$ both on the original exact $A$ and
at $y_\ell$, such that $\beta$ lies outside the union of the last Gamma
pairs of $P.\mathit{counterexample}$ and the base-extra missing set
of $P$ (computed using its own base core and terminal ladder).
\end{helper}
\begin{proof}
Apply \noderef{HierarchyCounterexamplePenultimateElementary} to
$\mathit{ctx.proper}$, $\mathit{ctx.oneEdgeUncolored}$,
$\mathit{ctx.lowerRung}$, $\mathit{ctx.fewerDividers}$,
$\mathit{ctx.lexicographicMinimality}$, and
$\mathit{ctx.largeMissingSet}$, with the ambient data in the statement.
It makes $B_{\rm pen}$ elementary under $\varphi$.
The fields of \noderef{TashkinovHierarchy} and the entrywise definition
of \noderef{TreeSequenceInitialSegment} give $A\preccurlyeq T$.
Since $\ell=|T|-|A|\ge3$, the exact prefix $B_{\rm pen}$ has length
$|T|-1$ and contains $A$. The base core is below $A$: at $q=0$
it equals $A$, and at $q>0$ use the core-inside-first field followed
by the nested stages. Also $T_n\preccurlyeq\mathcal B.\mathit{core}$:
in PE this is the left-segment field of the ordered union in
\noderef{HierarchyBase}, and in non-PE the core equals $T_n$.
Consequently $T_n$ and the base core both lie in $B_{\rm pen}$.

By \noderef{GoodHierarchyReservedDisjoint}, every last Gamma pair lies
in $\overline\varphi(A)$ and $\Gamma\cap D=\varnothing$.
Expanding \noderef{HierarchyStoredColorsAt}, one also has
$D=D_n\setminus\overline\varphi(A)$, where
$D_n=\noderef{TashkinovStoredColors}(\mathcal E.\mathit{series.data},n)$.
In particular no stored color is missing on $A$.

By \noderef{HierarchyStoredColorsCardinality}, $|D|\le n$.
The two-colors clause of \noderef{GoodHierarchy} gives
$|\Gamma|\le\sum_{\rho\in D}|C.\mathit{good}.\Gamma^q_\rho|
=2|D|\le2n$. A color cannot be missing both on $T_n$ and at a
vertex in $V(\mathcal B.\mathit{core})\setminus V(T_n)$:
these would be two distinct vertices of the elementary $B_{\rm pen}$.
Hence $\overline\varphi(T_n)\cap M_0=\varnothing$. The large-set
field of the context therefore yields
\[
 |\overline\varphi(T_n)\setminus(\Gamma\cup M_0)|
 \ge (2n+11)-2n=11.
\]
Choose $\beta$ from this set. It is missing on $T_n\preccurlyeq A$,
outside $\Gamma\cup M_0$, and different from $\alpha\in M_0$.
Both affected colors are missing on $A$ and consequently on
$B_{\rm pen}$; neither belongs to $D$.

Apply \noderef{ChangeLemma} in the order $(\beta,\alpha)$, using the
six context projections already listed and the derived elementary
prefix. Its conditional ordered-unused premise is vacuous because both
colors are missing on $A$. If $q>0$, then $A=T_{n,q}$ and $\beta$
is missing there, so the first conclusion identifies the components
through the first $\beta$- and $\alpha$-missing vertices. If $q=0$,
the chosen $\beta$ is missing on $T_n$, and the second conclusion
gives the same identification. This orientation does not require
$\alpha\notin D_n$.

The two first missing witnesses belong to $A\subseteq B_{\rm pen}$
and are distinct from $y_\ell$. They need not be distinct from each
other. Apply \noderef{P2905_T2_ElementaryTerminalComponent} to these
witnesses, the elementary vertex set of $B_{\rm pen}$, and $w=y_\ell$;
the color missing at $w$ is $\alpha$, the second color in this ordering.
It gives disjointness of the terminal component, its possibly singleton
path certificate, and equality of all missing sets on $B_{\rm pen}$.
Put $\sigma=\noderef{TwoColorComponentSwap}(G,\varphi,\beta,\alpha,y_\ell)$.
The missing-color exchange clause of
\noderef{TwoColorComponentSwapMissingBehavior} makes $\beta$ missing
at $y_\ell$ and the prefix equality retains its witness on $A$.
The entire missing set on $B_{\rm pen}$ is now fixed. Therefore every
totalized \noderef{HierarchyFirstMissingPosition} on $T$ is fixed,
in particular that of every $\rho\in D$. Indeed, if a color has a
witness before the last vertex, equality at every preceding vertex
preserves exactly its nonempty set of early candidates and its least
candidate. If it has no early witness, neither coloring has one,
and both totalized positions are the last index, whether or not a
terminal witness exists. This checks the fallback as well as the
actual-witness case.

Invoke \noderef{StableChangeCasesOneTwoStarredTail} with first color
$\beta$, second color $\alpha$, root the terminal vertex-node,
and this literal swap. The required two old missing colors on
$B_{\rm pen}$, path certificate, component disjointness, and Cases 1/2
alternative were all verified above. Its starred-missing premise is
the retained missingness of both colors on $A$; its Gamma-outside
premise follows from the choices; its pointwise penultimate missingness
and universally quantified stored-first premise were just checked.
This is the starred-tail reconstruction, so no assertion that all
post-$A$ displayed edges retain their values is needed.
We describe the identity presentation used below. For every slot of
the certified \noderef{TashkinovSeries}, properness and exact
one-uncoloredness follow by cases on its
\noderef{TashkinovSeriesStep}: the initial case uses the two coloring
clauses of its \noderef{KTriple}; SE uses
\noderef{MaximumDefectiveChoice}'s selected-coloring clauses; PE uses
the explicit current-coloring clauses; RE copies the previous coloring.
Choose the same input, series, and ETT for the reflexive
\noderef{CorrespondingExtendedTashkinovTree}; use reflexive rung,
presentation, ladder, stored-set, connecting-edge, and kind equalities,
the original generating-coloring equality, the just verified slotwise
coloring laws, and reflexive \noderef{StableColoring} at each positive
slot. Define $P_0$ by coloring $\varphi$, that certificate, maximum
$C.\mathit{maximum}$, base $\mathcal B$, counterexample $C$, reflexive
\noderef{SameGoodHierarchy}, and the supplied context $\mathit{ctx}$.
These are all seven fields of
\noderef{HierarchyReplacementPresentation}.
Retain the provider's full corresponding certificate $\mathcal C$,
its maximum witness, base $\mathcal B'$, hierarchy $K$, good witness
$g$, and $\noderef{SameGoodHierarchy}(C.\mathit{good},g)$.
Construct a \noderef{HierarchyCounterexample} $C'$ on $\mathcal B'$:
its positive-rung field is $C.\mathit{positive}$ transported by
$\mathcal C.\mathit{rungAgreement}$, its hierarchy and goodness are
$K,g$, and its maximum is the provider's maximum witness.
The two stated missing witnesses, one in the strict prefix
$B_{\rm pen}$ and one at $y_\ell$, are distinct by injectivity of
\noderef{TreeSequence}'s vertex display. Their common missing color
contradicts the disjointness clause of
\noderef{ColoredVertexSetElementary}, proving the non-elementarity
field. Apply \noderef{HierarchyReplacementPresentationTrans} to $P_0$,
$\sigma$, $\mathcal C$, $\mathcal B'$, $C'$, and the supplied same-good
witness. It gives an original-relative presentation $P$ with coloring
exactly $\sigma$; in particular its minimum context is supplied by the
principal, not assumed.

For any such $P$, its own same-good field identifies its base core,
last stored set, and every last Gamma pair with the originals after the
divider-count cast. The rung and ladder agreement in its corresponding
certificate identifies its actual terminal ladder tree with $T_n$:
use the last positive slot, which exists since $C.\mathit{positive}$
gives $n>0$, and its series-presentation equality. Missing-set equality
on the old base core and $T_n$ then identifies its base-extra missing
set with $M_0$. Thus the final forbidden set in the statement is the
original $\Gamma\cup M_0$.

Here the non-elementarity color is $\beta$, with its retained witness
on $A\subseteq B_{\rm pen}$. Since the original choice of $\beta$
excluded $\Gamma\cup M_0$, the final identification proves precisely
the forbidden-set exclusion for $P,\beta$. These are all the asserted
properties.
\end{proof}
